\documentclass[12pt, reqno]{amsart}
\usepackage{latexsym, amssymb, amscd, xypic, amsthm, setspace}
\input xy
\xyoption{all}

%%: Font = times
\usepackage{mathptmx} % rm & math
\usepackage[scaled=0.90]{helvet} % ss
\usepackage{courier} % tt
\normalfont
\usepackage[T1]{fontenc}
\pagestyle{plain}


\newtheorem{theorem}{Theorem}[subsection]
\newtheorem{prop}[theorem]{Proposition}
\newtheorem{lemma}[theorem]{Lemma}
\newtheorem{corollary}[theorem]{Corollary}
\newtheorem{claim}[theorem]{Claim}

\theoremstyle{definition}
\newtheorem{de}[theorem]{Definition}
\newtheorem{notation}{Notation}
\newtheorem{example}[theorem]{Example}
\newtheorem{examples}{Examples}
\newtheorem{remark}[theorem]{Remark}
\newtheorem{remarks}{Remarks}
\newtheorem{question}{Question}


\newtheorem*{nthm}{Theorem}
\newtheorem*{nlem}{Lemma}
\newtheorem*{nprop}{Proposition}
\newtheorem*{ncor}{Corollary}
\newtheorem*{nclaim}{Claim}
\newtheorem*{nde}{Definition}
\newtheorem*{nex}{Example}
\usepackage[most]{tcolorbox}
\definecolor{solblue}{RGB}{232,243,252}
\definecolor{solframe}{RGB}{120,170,215}
\newenvironment{solution}
  {\par\smallskip\begin{tcolorbox}[breakable, enhanced, colback=solblue, colframe=solframe,
     boxrule=0.5pt, arc=2pt, left=6pt, right=6pt, top=4pt, bottom=4pt]\noindent\textbf{Solution.}}
  {\hfill$\square$\end{tcolorbox}\par\medskip}

\newtheoremstyle{test}{3pt}{3pt}{\itshape}{\parindent}{\bfseries}{---}{  }{\thmnote{#3}} \theoremstyle{test}
\newtheorem*{mythm}{}


%: paragraphs

\newcommand{\point}{\vspace{3mm}\par \noindent (\refstepcounter{subsection}{\bf \thesubsection) } }

\newcommand{\tpoint}[1]{\vspace{3mm}\par \noindent \refstepcounter{subsubsection}{\bf \thesubsubsection.}
  {\em #1. ---} }
\newcommand{\epoint}[1]{\vspace{3mm}\par \noindent \refstepcounter{subsection}{\bf \thesubsection.}
  {\em #1.} }
\newcommand{\spoint}{\vspace{3mm}\par \noindent \refstepcounter{subsection}{\bf \thesubsection.} }

\newcommand{\bpoint}[1]{\vspace{3mm}\par \noindent \refstepcounter{subsection}{\bf \thesubsection.}
  {\bf #1.} }



%: numbering and margins

\numberwithin{equation}{subsection}

\setlength{\oddsidemargin}{0cm} \setlength{\evensidemargin}{0cm}
\setlength{\marginparwidth}{0in}
\setlength{\marginparsep}{0in}
\setlength{\marginparpush}{0in}
\setlength{\topmargin}{0in}
\setlength{\headheight}{0pt}
\setlength{\headsep}{0pt}
\setlength{\footskip}{.3in}
\setlength{\textheight}{9.2in}
\setlength{\textwidth}{6.5in}
\setlength{\parskip}{0.25pt}
\setlength{\parindent}{0.25in}

%: math shortcuts


\DeclareMathOperator{\inv}{inv}
\DeclareMathOperator{\out}{Out}
\DeclareMathOperator{\gal}{Gal}
\renewcommand{\mod}{\mathrm{\, mod \, }}
\renewcommand{\d}[1]{{}^{#1}}
\renewcommand{\b}[1]{\mathbf{{#1}}}
\newcommand{\C}{\mathbb{C}}
\newcommand{\be}[1]{\begin{eqnarray} \label{#1}}
\newcommand{\ee}{\end{eqnarray}}
\newcommand{\re}{\mathsf{Re}}
\newcommand{\im}{\mathsf{Im}}
\newcommand{\rr}{\rightarrow}

\newcommand{\n}{\mathbb{N}}
\newcommand{\zee}{\mathbb{Z}}
\newcommand{\R}{\mathbb{R}}
\newcommand{\pr}{\mathbb{P}}
\renewcommand{\gcd}[2]{\mathsf{gcd}(#1, #2)}
\newcommand{\ov}[1]{\overline{#1}}
\renewcommand{\Re}{\mathrm{Re}}
\renewcommand{\Im}{\mathrm{Im}}



\title{Problem Set 2: Solutions \\ MATH 411: Honors Complex Variables, Fall 2026}
\date{last compiled: \today}

\numberwithin{equation}{section}
\refstepcounter{section}

\begin{document}						%%
\maketitle


\begin{enumerate}

%%%%%%%%%%%%%%%%%%%%%%%%%%%%%%%%%%%%%%%%%%%%%%%%%%%%%%%%%%%%%%%%%%%%%%%
\item Give the terms of order $\leq 3$ in the power series (using the expressions for $\cos(T)$, $\sin(T)$, and $e^T$ given in lecture):
\begin{enumerate}
\item $\dfrac{\sin T}{\cos T}$;
\item $\dfrac{e^T}{\sin T}$.
\end{enumerate}

\begin{solution}
In class we defined
\begin{align*}
\cos T&=1-\frac{T^2}{2!}+\frac{T^4}{4!}-\cdots=\sum_{n=0}^{\infty}(-1)^n\frac{T^{2n}}{(2n)!},\\
\sin T&=T-\frac{T^3}{3!}+\frac{T^5}{5!}-\cdots=\sum_{n=0}^{\infty}(-1)^n\frac{T^{2n+1}}{(2n+1)!},\\
e^T&=1+T+\frac{T^2}{2!}+\frac{T^3}{3!}+\cdots=\sum_{n=0}^{\infty}\frac{T^n}{n!}.
\end{align*}

\medskip
\noindent (a) Write
\[
\frac{1}{\cos T}=b_0+b_1T+b_2T^2+\cdots,
\]
so that
\[
\bigl(b_0+b_1T+b_2T^2+\cdots\bigr)\Bigl(1-\frac{T^2}{2!}+\frac{T^4}{4!}-\cdots\Bigr)=1 .
\]
Comparing coefficients:
\begin{align*}
b_0&=1 &&\text{(coefficient of $T^0$)},\\
b_1&=0 &&\text{(coefficient of $T$)},\\
b_2-\frac{b_0}{2}&=0,\ \text{so } b_2=\tfrac12 &&\text{(coefficient of $T^2$)}.
\end{align*}
Hence
\[
\frac{1}{\cos T}=1+\frac12T^2+\cdots
\]
and
\begin{align*}
\frac{\sin T}{\cos T}&=\Bigl(1+\frac12T^2+\cdots\Bigr)\Bigl(T-\frac{T^3}{3!}+\frac{T^5}{5!}-\cdots\Bigr)\\
&=T+\Bigl(\frac12-\frac16\Bigr)T^3+\cdots=\boxed{T+\frac{T^3}{3}}+\cdots .
\end{align*}

\medskip
\noindent (b) Write
\[
\sin T=T\Bigl(1-\frac{T^2}{3!}+\frac{T^4}{5!}-\cdots\Bigr)=T\bigl(1-h(T)\bigr),
\]
where
\[
h(T)=\frac{T^2}{3!}-\frac{T^4}{5!}+\cdots .
\]
Since $h$ has no constant term,
\begin{align*}
\frac{1}{\sin T}&=\frac1T\bigl(1+h(T)+h(T)^2+\cdots\bigr)
=\frac1T\Bigl(1+\frac{T^2}{3!}+\Bigl(\frac1{36}-\frac1{120}\Bigr)T^4+\cdots\Bigr)\\
&=\frac1T\Bigl(1+\frac{T^2}{3!}+\frac{7T^4}{360}+\cdots\Bigr).
\end{align*}
So
\begin{align*}
\frac{e^T}{\sin T}&=\frac1T\Bigl(1+T+\frac{T^2}{2!}+\frac{T^3}{3!}+\frac{T^4}{4!}+\cdots\Bigr)\\
&\qquad\cdot\Bigl(1+\frac{T^2}{3!}+\frac{7T^4}{360}+\cdots\Bigr)\\
&=\frac1T\Bigl(1+T+\frac23T^2+\frac13T^3+\frac{13}{90}T^4+\cdots\Bigr)\\
&=\boxed{\frac1T+1+\frac{2T}{3}+\frac{T^2}{3}+\frac{13T^3}{90}}+\cdots .
\end{align*}
(For the $T^4$ coefficient inside the parentheses: $\frac1{24}+\frac1{12}+\frac7{360}=\frac{52}{360}=\frac{13}{90}$.)
\end{solution}

%%%%%%%%%%%%%%%%%%%%%%%%%%%%%%%%%%%%%%%%%%%%%%%%%%%%%%%%%%%%%%%%%%%%%%%
\item Given complex numbers $a_0,a_1,u_1,u_2$, define $a_n$ for $n\geq 2$ by
\[
a_n=u_1a_{n-1}+u_2a_{n-2}.
\]
If
\[
T^2-u_1T-u_2=(T-\alpha)(T-\alpha')
\]
with $\alpha\neq\alpha'$, show that
\[
a_n=A\alpha^n+B(\alpha')^n
\]
for suitable $A,B$. Find $A,B$ in terms of $a_0,a_1,\alpha,\alpha'$. Consider the power series
\[
F(T)=\sum_{n=0}^{\infty}a_nT^n.
\]
Show that it represents a rational function $f(T)/g(T)$, where $f(T)$ and $g(T)$ are polynomials with $g(T)\neq 0$. Give the partial fraction decomposition of this rational function.

\begin{solution}
If $T^2-u_1T-u_2=(T-\alpha)(T-\alpha')$, then we first note that
\begin{equation}\label{u:alpha}
u_1=\alpha+\alpha'\quad\text{and}\quad u_2=-\alpha\alpha' .
\end{equation}
Now, if $a_n=A\alpha^n+B(\alpha')^n$, we must have in particular that
\begin{equation}\label{a12}
a_1=A\alpha+B\alpha'\quad\text{and}\quad a_0=A+B.
\end{equation}
These can be solved for $A,B$ to yield
\begin{equation}\label{A:B}
A=\frac{a_1-a_0\alpha'}{\alpha-\alpha'}\quad\text{and}\quad B=a_0-A=\frac{a_0\alpha-a_1}{\alpha-\alpha'} .
\end{equation}
Note that the ratio involved in the definition of $A$ is well-defined as we assumed $\alpha\neq\alpha'$. So let us now
\emph{define} $A$ and $B$ as in \eqref{A:B} and verify that with these definitions we must have
\begin{equation}\label{an}
a_n=A\alpha^n+B(\alpha')^n\quad\text{for any } n\geq0.
\end{equation}
Indeed, for $n=0,1$ this follows immediately from the way we have defined $A$ and $B$, so let us assume the formula
\eqref{an} holds for $a_n$ when $n=0,\dots,m$ with $m\geq1$. We argue that it must also hold for $m+1$, i.e.\
$a_{m+1}=A\alpha^{m+1}+B(\alpha')^{m+1}$. Indeed, by hypothesis, we have
\begin{align*}
a_{m+1}&=u_1a_m+u_2a_{m-1}=u_1\bigl(A\alpha^m+B(\alpha')^m\bigr)+u_2\bigl(A\alpha^{m-1}+B(\alpha')^{m-1}\bigr)\\
&=A\alpha^{m-1}(u_1\alpha+u_2)+B(\alpha')^{m-1}(u_1\alpha'+u_2).
\end{align*}
Using \eqref{u:alpha} (equivalently, since $\alpha,\alpha'$ are roots of $T^2-u_1T-u_2$), we have $u_1\alpha+u_2=\alpha^2$
and $u_1\alpha'+u_2=(\alpha')^2$, so the induction is proven.

\medskip
Defining $F(T)$ as in the problem, we find
\begin{align*}
u_1TF(T)&=\sum_{n\geq1}u_1a_{n-1}T^n=u_1a_0T+u_1a_1T^2+\cdots,\\
u_2T^2F(T)&=\sum_{n\geq2}u_2a_{n-2}T^n=u_2a_0T^2+u_2a_1T^3+\cdots,
\end{align*}
and therefore, using the recursion for the coefficients of $T^n$ with $n\geq2$, we have (formally)
\[
F(T)-a_0-a_1T=u_1TF(T)-u_1a_0T+u_2T^2F(T).
\]
Solving for $F(T)$ we find
\begin{equation}\label{sol:F}
F(T)=\frac{a_0+a_1T-u_1a_0T}{1-u_1T-u_2T^2},
\end{equation}
so $F=f/g$ with $f(T)=a_0+(a_1-u_1a_0)T$ and $g(T)=1-u_1T-u_2T^2$; note $g\neq0$ since its constant term is $1$.
From \eqref{u:alpha} we find that
\[
(1-\alpha T)(1-\alpha'T)=1-(\alpha+\alpha')T+\alpha\alpha'T^2=1-u_1T-u_2T^2,
\]
so that
\[
F(T)=\frac{a_0+a_1T-u_1a_0T}{(1-\alpha T)(1-\alpha'T)}=\boxed{\frac{A}{1-\alpha T}+\frac{B}{1-\alpha'T}},
\]
with $A,B$ as in \eqref{A:B}. This last equality follows immediately from clearing denominators and using \eqref{u:alpha}
and \eqref{a12}: indeed $A(1-\alpha'T)+B(1-\alpha T)=(A+B)-(A\alpha'+B\alpha)T$, and
$A\alpha'+B\alpha=(A+B)(\alpha+\alpha')-(A\alpha+B\alpha')=u_1a_0-a_1$.
(Alternatively, expand each term as a geometric series, $\frac{A}{1-\alpha T}=\sum A\alpha^nT^n$, and compare with \eqref{an}.)
\end{solution}

%%%%%%%%%%%%%%%%%%%%%%%%%%%%%%%%%%%%%%%%%%%%%%%%%%%%%%%%%%%%%%%%%%%%%%%
\item Determine the radius of convergence of the following series:
\begin{enumerate}
\item $\displaystyle\sum n^2z^n$;
\item $\displaystyle\sum \frac{n!}{n}z^n$;
\item If $f=\sum a_nz^n$ has radius of convergence $r>0$, show that $\sum n^da_nz^n$ has the same radius of convergence for any positive integer $d$.
\end{enumerate}

\begin{solution}
\noindent (a) It is $\boxed{1}$. We use Hadamard's criterion and compute
\[
\limsup_{n\to\infty}|n^2|^{1/n}=\limsup_{n\to\infty}n^{2/n}=\limsup_{n\to\infty}e^{\frac{2\log n}{n}}=1,
\]
since $\lim_{n\to\infty}\frac{2\log n}{n}=0$.

\medskip
\noindent (b) First note that
\[
\lim_{n\to\infty}(1/n)^{1/n}=1 .
\]
Also, as we observed in class, we have the Stirling estimate $n!=n^ne^{-n}u_n$ where $\lim u_n^{1/n}=1$. Thus
$(n!)^{1/n}=n\,e^{-1}u_n^{1/n}$, and
\[
\limsup\,(n!)^{1/n}(1/n)^{1/n}=\limsup\,n\,e^{-1}u_n^{1/n}(1/n)^{1/n}=\infty,
\]
so the radius of convergence is $\boxed{0}$.

\medskip
\noindent (c) Again, we use Hadamard's criterion to compute $\limsup_{n\to\infty}|n^da_n|^{1/n}$. First,
\[
\lim_{n\to\infty}n^{d/n}=\lim_{n\to\infty}e^{\frac{d\log n}{n}}=1 .
\]
Since $|n^da_n|^{1/n}=n^{d/n}|a_n|^{1/n}$, and $n^{d/n}$ is a sequence of positive numbers with limit $1$, the remarks
before Theorem 2.6 in Lang (p.~55) (if $\lim t_n=t$ exists and $t\neq0,\infty$, then $\limsup t_ns_n=t\limsup s_n$)
give
\[
\limsup_{n\to\infty}|n^da_n|^{1/n}=\limsup_{n\to\infty}|a_n|^{1/n}=1/r .
\]
Hence $\sum n^da_nz^n$ also has radius of convergence $r$.
\end{solution}

%%%%%%%%%%%%%%%%%%%%%%%%%%%%%%%%%%%%%%%%%%%%%%%%%%%%%%%%%%%%%%%%%%%%%%%
\item Prove the following:
\begin{enumerate}
\item The power series $\sum nz^n$ has radius of convergence $1$ but does not converge at any point of the unit circle.
\item The power series $\sum \frac{z^n}{n^2}$ has radius of convergence $1$ and converges at every point of the unit circle.
\item The power series $\sum \frac{z^n}{n}$ converges at every point of the unit circle except $z=1$.
\end{enumerate}

\begin{solution}
\noindent (a) Since $\limsup n^{1/n}=1$, the assertion about the radius of convergence being equal to $1$ follows from
Hadamard's theorem. On the other hand, if $|z|=1$, then the terms $nz^n$ of the series have norm $n$, which approaches
infinity, so the series clearly cannot converge for any such $z$.

\medskip
\noindent (b) Since $\limsup\frac{1}{n^{2/n}}=1$, the radius of convergence is $1$. If $|z|=1$, then
\[
\sum\Bigl|\frac{z^n}{n^2}\Bigr|=\sum\frac1{n^2},
\]
and the latter converges by comparison with (a Riemann sum for) the convergent integral $\int_1^\infty\frac{1}{x^2}\,dx$, for
example. So the series converges absolutely at every point of the unit circle.

\medskip
\noindent (c) Again, from Hadamard we see the radius of convergence is $1$. For $z=1$, this is the harmonic series, whose
divergence is well known. So let us assume that $z\neq1$ but $z$ lies on the unit circle.

We use \emph{summation by parts}: if $\{a_n\}_{n=1}^N$ and $\{b_n\}_{n=1}^N$ are finite sequences of complex numbers and
$B_k=\sum_{n=1}^kb_n$, then
\begin{equation}\label{sbp}
\sum_{n=1}^Na_nb_n=a_NB_N-\sum_{n=1}^{N-1}(a_{n+1}-a_n)B_n .
\end{equation}
(Proof by induction on $N$: for $N=1$ both sides are $a_1b_1$. If $L_N$, $R_N$ denote the left and right hand sides, then
$L_{N+1}-L_N=a_{N+1}b_{N+1}$, while
$R_{N+1}-R_N=a_{N+1}B_{N+1}-a_NB_N-(a_{N+1}-a_N)B_N=a_{N+1}(B_{N+1}-B_N)=a_{N+1}b_{N+1}$.)

Apply \eqref{sbp} to $\sum_{n=1}^N\frac1n z^n$ with $a_n=1/n$ and $b_n=z^n$. As $z\neq1$ we may write
$B_n=\frac{z-z^{n+1}}{1-z}$, and since $a_{n+1}-a_n=-\frac{1}{n(n+1)}$, we obtain
\[
\sum_{n=1}^N\frac1n z^n=\frac{1}{1-z}\cdot\frac{z-z^{N+1}}{N}+\frac{1}{1-z}\sum_{n=1}^{N-1}\frac{z-z^{n+1}}{n(n+1)} .
\]
Now if $z$ is on the unit circle we have $|z-z^k|\leq2$ for any integer $k$. Hence the first term is bounded in norm by
$\frac{2}{|1-z|N}$, which tends to $0$ as $N\to\infty$. In the second term, the $n$-th summand of the sum is bounded in norm by
$\frac{2}{n(n+1)}$, and $\sum\frac{2}{n(n+1)}$ converges (it telescopes: $\frac1{n(n+1)}=\frac1n-\frac1{n+1}$). So the series
$\sum_{n\geq1}\frac{z-z^{n+1}}{n(n+1)}$ converges absolutely. Letting $N\to\infty$, the partial sums
$\sum_{n=1}^N\frac{z^n}{n}$ converge, i.e.\ $\sum\frac{z^n}{n}$ converges at every $z\neq1$ with $|z|=1$.
\end{solution}

%%%%%%%%%%%%%%%%%%%%%%%%%%%%%%%%%%%%%%%%%%%%%%%%%%%%%%%%%%%%%%%%%%%%%%%
\item Find the terms of order $\leq 3$ in the power series expansion of
\[
f(z)=\frac{z-2}{(z+3)(z+2)}
\]
at $z=1$.

\begin{solution}
We begin by noting that if $|z/a|<1$ we have
\[
\frac{1}{a+z}=\frac1a\sum_{n=0}^{\infty}(-1)^n\Bigl(\frac za\Bigr)^n .
\]
Note that
\[
z-2=-1+(z-1),\qquad z+3=4+(z-1),\qquad z+2=3+(z-1).
\]
Hence
\begin{align*}
\frac{1}{z+3}&=\frac{1}{4+(z-1)}=\frac14\sum_{n=0}^{\infty}(-1)^n\Bigl(\frac{z-1}{4}\Bigr)^n,\\
\frac{1}{z+2}&=\frac{1}{3+(z-1)}=\frac13\sum_{n=0}^{\infty}(-1)^n\Bigl(\frac{z-1}{3}\Bigr)^n,
\end{align*}
and so, for $|z-1|<3$,
\begin{align*}
\frac{z-2}{(z+3)(z+2)}&=\frac1{12}\bigl(-1+(z-1)\bigr)\sum_{n=0}^{\infty}(-1)^n\Bigl(\frac{z-1}{4}\Bigr)^n\sum_{n=0}^{\infty}(-1)^n\Bigl(\frac{z-1}{3}\Bigr)^n\\
&=\frac1{12}\bigl(-1+(z-1)\bigr)\\
&\qquad\cdot\Bigl(1-\frac{7}{12}(z-1)+\frac{37}{144}(z-1)^2-\frac{175}{1728}(z-1)^3+\cdots\Bigr)\\
&=\boxed{-\frac1{12}+\frac{19}{144}(z-1)-\frac{121}{1728}(z-1)^2+\frac{619}{20736}(z-1)^3}+\cdots
\end{align*}
in a neighbourhood of $z=1$. (Here the coefficient of $(z-1)^n$ in the product of the two geometric series is
$(-1)^n\sum_{k=0}^n4^{-k}3^{-(n-k)}$; e.g.\ for $n=2$ it is $\frac1{16}+\frac1{12}+\frac19=\frac{37}{144}$.)
\end{solution}

\end{enumerate}

\end{document}
